The bug report points out a real gap: one must justify first that the exposed face \([t_i,t_{i+1}]\) is parametrized by \(a_i\) with a **nonnegative** coefficient. Once that is inserted, the rest of the argument goes through.

Here is a corrected proof.

---

### Proof

Fix \(i\), and abbreviate
\[
a:=a_i=s_{i+1}-s_i,\qquad
x_i:=s_i+t_i,\qquad
x_{i+1}:=s_{i+1}+t_{i+1},
\]
\[
E:=[x_i,x_{i+1}],\qquad
F:=F_i=T\cap[t_i,t_{i+1}].
\]

From Subproblems \(1\)–\(3\), we may use the following facts.

1. \(E\) is the exposed face of \(P+Q\) corresponding to the edge \([s_i,s_{i+1}]\), and
   \[
   E=[s_i,s_{i+1}]+[t_i,t_{i+1}].
   \]
2. \(E\) is an edge of \(P+Q\) parallel to \(a\).
3. Every interior point of \(E\) cancels:
   \[
   \sigma(p)=0\qquad\text{for all }p\in\operatorname{relint}(E).
   \]
4. If \(p\in E\) and \(p=t+s\) with \(t\in T\) and \(s\in S\), then
   \[
   t\in[t_i,t_{i+1}],\qquad s\in[s_i,s_{i+1}].
   \]

Also, since \(S\) is in convex position, the edge \([s_i,s_{i+1}]\) contains no point of \(S\) other than its endpoints:
\[
S\cap [s_i,s_{i+1}]=\{s_i,s_{i+1}\}.
\tag{1}
\]

---

### Step 1: justify the parametrization of \([t_i,t_{i+1}]\)

Because \(E\) is parallel to \(a\), the vector
\[
x_{i+1}-x_i=(s_{i+1}-s_i)+(t_{i+1}-t_i)=a+(t_{i+1}-t_i)
\]
is parallel to \(a\). Hence \(t_{i+1}-t_i\) is also parallel to \(a\), so there exists \(M\in\mathbb{R}\) such that
\[
t_{i+1}=t_i+Ma.
\tag{2}
\]

We now show \(M\ge 0\).

If \(M<0\), then
\[
[t_i,t_{i+1}]=t_i+[M,0]\,a.
\]
Using the face identity from Subproblem \(3\),
\[
E=[s_i,s_{i+1}]+[t_i,t_{i+1}]
  =x_i+[0,1]a+[M,0]a
  =x_i+[M,1]a.
\]
But then \(x_i=x_i+0\cdot a\) is an interior point of this segment, not an endpoint, contradicting \(E=[x_i,x_{i+1}]\).

Therefore
\[
M\ge 0.
\]
So (2) becomes
\[
[t_i,t_{i+1}]
=\{\,t_i+\lambda a:0\le \lambda\le M\,\},
\qquad
E=x_i+[0,M+1]a.
\tag{3}
\]

Define
\[
\Lambda:=\{\lambda\in[0,M]: t_i+\lambda a\in F\}.
\]
Since \(t_i,t_{i+1}\in F\), we have
\[
0,M\in\Lambda.
\tag{4}
\]

---

### Step 2: successor and predecessor relations

Let \(\lambda\in\Lambda\) with \(\lambda<M\), and set
\[
t:=t_i+\lambda a,\qquad p:=s_{i+1}+t.
\]
By (3),
\[
p=s_{i+1}+t_i+\lambda a=x_i+(\lambda+1)a.
\]
Since \(0<\lambda+1<M+1\), we have \(p\in\operatorname{relint}(E)\). Hence \(\sigma(p)=0\).

One contribution to \(\sigma(p)\) comes from \(t\), because
\[
p-t=s_{i+1}.
\]
Any other contribution must come from some \(t'\in T\) with \(p-t'\in S\). By Fact 4, \(t'\in[t_i,t_{i+1}]\) and \(p-t'\in[s_i,s_{i+1}]\). By (1), the only possibilities for \(p-t'\) are \(s_i\) and \(s_{i+1}\). The case \(s_{i+1}\) gives \(t'=t\); the case \(s_i\) gives
\[
t'=p-s_i=t+a.
\]
So the only possible contributors are \(t\) and \(t+a\). Since \(\sigma(p)=0\) and \(\sigma(t)\in\{\pm1\}\), the second one must actually occur and must cancel the first:
\[
t+a\in F,\qquad \sigma(t+a)=-\sigma(t).
\tag{5}
\]

Similarly, let \(\lambda\in\Lambda\) with \(\lambda>0\), and set
\[
t:=t_i+\lambda a,\qquad p:=s_i+t=x_i+\lambda a.
\]
Again \(p\in\operatorname{relint}(E)\), so \(\sigma(p)=0\). As above, the only possible contributors are \(t\) and
\[
p-s_{i+1}=t-a.
\]
Hence
\[
t-a\in F,\qquad \sigma(t-a)=-\sigma(t).
\tag{6}
\]

In terms of \(\Lambda\), (5) and (6) say:
\[
\lambda\in\Lambda,\ \lambda<M \implies \lambda+1\in\Lambda,
\tag{7}
\]
\[
\lambda\in\Lambda,\ \lambda>0 \implies \lambda-1\in\Lambda.
\tag{8}
\]

---

### Step 3: there are no fractional parameters

We claim that
\[
\Lambda\subset\mathbb{Z}.
\]

Assume not, and choose \(\lambda\in\Lambda\setminus\mathbb{Z}\). Write
\[
r:=\lfloor\lambda\rfloor,\qquad \lambda_0:=\lambda-r.
\]
Then
\[
0<\lambda_0<1.
\]
Applying (8) repeatedly \(r\) times yields
\[
\lambda_0\in\Lambda.
\]

Now let
\[
t:=t_i+\lambda_0 a,\qquad p:=s_i+t=x_i+\lambda_0 a.
\]
Because \(0<\lambda_0<1\), we have \(p\in\operatorname{relint}(E)\), so \(\sigma(p)=0\).

But the same two-contributor analysis as above shows that the only possible contributors to \(\sigma(p)\) are \(t\) and \(t-a\). Here
\[
t-a=t_i+(\lambda_0-1)a,
\]
and since \(\lambda_0-1<0\), this point is not on \([t_i,t_{i+1}]\), hence not in \(F\). Therefore \(p\) has exactly one contribution, namely \(t\), and so
\[
\sigma(p)=\sigma(t)\in\{\pm1\},
\]
contradicting \(\sigma(p)=0\).

Thus indeed
\[
\Lambda\subset\mathbb{Z}.
\tag{9}
\]

Since \(M\in\Lambda\) by (4), (9) implies
\[
M\in\mathbb{Z}_{\ge 0}.
\]
Set
\[
m_i:=M.
\]

---

### Step 4: the face is the full arithmetic progression and the signs alternate

From \(0\in\Lambda\) and the successor rule (7), we obtain
\[
1,2,\dots,m_i\in\Lambda.
\]
Together with \(\Lambda\subset[0,m_i]\cap\mathbb{Z}\), this gives
\[
\Lambda=\{0,1,\dots,m_i\}.
\]
Therefore
\[
F_i=\{\,t_i+k\,a_i:0\le k\le m_i\,\},
\qquad
t_{i+1}=t_i+m_i a_i.
\]

Finally, (5) gives for \(0\le k<m_i\),
\[
\sigma(t_i+(k+1)a_i)=-\sigma(t_i+k a_i).
\]
Induction on \(k\) yields
\[
\sigma(t_i+k a_i)=(-1)^k\,\sigma(t_i)
\qquad(0\le k\le m_i).
\]

So every \(F_i\) is a full arithmetic progression in direction \(a_i\), and the signs alternate along it. \(\square\)

---

If you want, I can also rewrite this in a shorter “official solution” style.